\section{Methodology}
\label{sec:methodology}

Section~\ref{sec:system} describes the system we target and how we account
for its energy. Section~\ref{sec:smart} presents SAGE-smart, and
Section~\ref{sec:protocols} compares the protocols that bring the state of the
clients into the decision.

% =====================================================================
\subsection{System Model}\label{sec:system}

We consider $N$ smartphones and one server that run $R$ synchronous rounds of
FL. In round $r$ the server selects a set $S_r$ of clients. Each selected
client downloads the global model, trains it on its $n_i$ local samples and
uploads the update. A round lasts as long as its slowest selected client.
Every phone that is not training, selected or not, is idle for the rest of the
round, as in~\cite{maciel2024escs}. Every term below is charged to the battery
of the device that spends it.

\textbf{Computation.}
Each phone belongs to one of three tiers, which are the three smartphones
profiled in~\cite{xu2022mandheling} (Table~\ref{tab:tiers}). A tier has a
training throughput $\Theta_i$ and a constant training power $P_i$, both
derived from the measured per-batch latency and energy of FP32 training on the
CPU. The battery capacities $Q_i$ are those of the same phones. Training $E$
epochs with batch size $B$ takes
\begin{equation}
  t_i(E,B,w) = E\left(\frac{3\,M\,n_i}{(1-w)\,\Theta_i}
             + \left\lfloor\frac{n_i}{B}\right\rfloor t_b\right),
  \label{eq:ttrain}
\end{equation}
where $M$ is the number of forward MACs per sample ($555.5\times10^{6}$ for
ResNet-18 on $32\times32$ images), the factor 3 accounts for the backward pass,
$t_b=57\,\mathrm{ms}$ is the fixed cost of a batch~\cite{xu2022mandheling}, and
$w\in[0,1)$ is the fraction of the phone taken by other applications. The
energy is $P_i\,t_i$. Since $P_i$ does not depend on $B$, a smaller batch costs
slightly more per epoch but performs more SGD steps per joule.

\begin{table}[t]
\centering
\caption{Device tiers: speed and power from~\cite{xu2022mandheling}, battery
capacities from the manufacturers.}
\label{tab:tiers}
\footnotesize
\setlength{\tabcolsep}{4pt}
\begin{tabular}{llrrr}
\hline
Tier & Phone (SoC chip) & $\Theta_i$ [GMAC/s] & $P_i$ [W] & $Q_i$ [mAh] \\
\hline
Low  & Redmi Note 9 Pro (SD750G) & 13.9 & 2.99 & 4820 \\
Mid  & Xiaomi Mi 10 (SD865)      & 26.6 & 4.00 & 4780 \\
High & Xiaomi Mi 11 Pro (SD888)  & 30.5 & 6.15 & 5000 \\
\hline
\end{tabular}
\end{table}

\textbf{Background workload and idle.}
The load $w$ of the other applications follows a Markov chain over three
states, $w\in\{0,\,0.10,\,0.30\}$, with one step per round. Three chains, with
stationary means 0.03, 0.12 and 0.21, are assigned to the phones in equal
shares and independently of the tier. The load slows the training down
(Eq.~\eqref{eq:ttrain}) and sets the idle power:
\begin{equation}
  P^{\mathrm{idle}}_i(w) = P_0 + w\,(P_i - P_0), \qquad P_0=0.1\,\mathrm{W}.
  \label{eq:idle}
\end{equation}
The phone knows its own $w$. The server sees it only through reports.

\textbf{Battery.}
The SoC of a phone drops by the charge drawn at its nominal voltage
($V=3.7\,\mathrm{V}$), and a phone at zero SoC is depleted and leaves the
federation. We consider three battery models, which define three simulated
worlds. The \emph{ideal} model draws $P\,t/V$. The \emph{Peukert} model scales
the drawn charge by $(I/I_R)^{n-1}$, with $I_R=0.2$C and $n=1.05$; we
obtained $n$ by least squares on the capacities of a commercial cell (2.60,
2.50 and 2.39~Ah at 0.2C, 0.5C and 1C)~\cite{doerffel2006peukert}. The
\emph{datasheet} model has no exponent. It interpolates the discharge curves
$V(\mathrm{SoC},C)$ and capacities $Q(C)$ of the same cell, normalized in
C-rate so that they apply to any capacity, and solves $P=V(\mathrm{SoC},I)\,I$
for the current. Its voltage drops with the charge, so at constant power the
discharge accelerates near the reserve, which is where the battery constraint
binds. \rev{[Cite the cell datasheet.]}

\textbf{Communication.}
The phones use LTE. A transfer of $s$~MB costs $e_{\mathrm{dl}}=0.98\,\mathrm{J/MB}$
in download and $e_{\mathrm{ul}}=2.96\,\mathrm{J/MB}$ in upload. These are the
energy-per-throughput fits of~\cite{narayanan2021variegated} at the median
Italian 4G throughputs of 25.1 and 10~Mbps~\cite{opensignal2018italy}.
Each isolated transfer also opens a radio session: a promotion
(1.21~W for 265~ms~\cite{huang2012lte,narayanan2021variegated})
and a tail (178~mW for 10.2~s~\cite{narayanan2021variegated}),
that is $E_{\mathrm{s}}\approx2.1\,\mathrm{J}$ whatever the payload. The tail costs
energy but does not delay the round. Exchanging ResNet-18 (44.9~MB each
way) costs
\begin{equation}
  E_{\mathrm{model}} = s\,(e_{\mathrm{dl}}+e_{\mathrm{ul}}) + 2E_{\mathrm{s}}
  \approx 181\,\mathrm{J},
  \label{eq:comm}
\end{equation}
with about 50~s of transfer. A control message (1~KB round
trip) costs a full session, $E_{\mathrm{s}}$, when it is isolated. It costs
almost nothing (about 2~mJ) when it shares the session of the model
download that follows. Each exchange delays the round by
$t_{\mathrm{c}}\approx0.33\,\mathrm{s}$ (promotion, RTT and transfer).

\textbf{User behavior.}
Initial SoCs are drawn from a normal distribution with mean 0.75 and standard
deviation 0.15, truncated to $[0.3,1]$. During a round of length $\Delta t$, a
phone is put on the charger with probability
$p = 1-\exp\!\left(-\lambda\,(1-\mathrm{SoC})\,\Delta t\right)$, with
$\lambda=0.126\,\mathrm{h}^{-1}$. This value gives one charging session per day at
the average SoC of 67\% observed in~\cite{ferreira2011battery}. A phone
charges at 1C until it is full, and it stays selectable meanwhile, training at
no cost for its battery. Real users rarely let the phone go below
30\%~\cite{ferreira2011battery}; ours do, so the scenario is pessimistic on
purpose.

% =====================================================================
\subsection{SAGE-smart}\label{sec:smart}

SAGE-smart splits the decision in two. The server chooses \emph{who}
participates, from what it knows. Each selected client chooses \emph{how
much} it trains, from what only it knows, at the moment training starts.
Fig.~\ref{alg:smart} summarizes one round.

\textbf{What the server knows.}
At enrollment, each client registers its static data: tier, capacity, battery
type, dataset size $n_i$ and label distribution, from which the server computes
the Jensen--Shannon divergence $D_i$ from the global label distribution, as in
SAGE~\cite{savoia2026sage}. Afterwards, the server only knows the last report
of each client: SoC, load $w$, charging status and the time $t_i^{\mathrm{rep}}$
at which they were read. Before ranking, it extrapolates the SoC with the idle
drain of Eq.~\eqref{eq:idle} through the battery model of the device:
\begin{equation}
  \widehat{\mathrm{SoC}}_i = \mathrm{SoC}_i^{\mathrm{rep}}
  - \Delta\mathrm{SoC}_i\!\left(P^{\mathrm{idle}}_i(w_i^{\mathrm{rep}}),\,
    t - t_i^{\mathrm{rep}}\right).
  \label{eq:estimate}
\end{equation}
A client that was charging at the last report is assumed unchanged. The
estimate does not see recharges or changes of load since the report: this is
the information that only the client has.

\textbf{Selection.}
Clients are ranked by
\begin{equation}
  s_i = a\,\widehat{\mathrm{SoC}}_i + b\,D_i
      + c\,\min\!\left(1, \eta_i/\eta_{\max}\right), \quad a+b=1,
  \label{eq:score}
\end{equation}
where $\eta_i$ is the number of rounds since client $i$ was last contacted. The
staleness term gives a chance to clients that the first two terms would never
pick: a client with low divergence, or one that drains slowly and stays
charged, would otherwise be excluded or selected at every round. We use $c=0.3$
and $\eta_{\max}=25$, and $a,b$ tuned per data distribution. As in SAGE, the
number of selected clients grows by one per round from two, up to 20\% of the
clients still alive. With non-negative scores and only cardinality
constraints, the selection ILP of SAGE reduces to taking the top-ranked
clients.

\textbf{Envelope.}
Instead of a fixed $(E,B)$, each selected client receives an envelope
$\mathcal{E}_i$: an epoch range $[E_{\mathrm{lo}},E_{\mathrm{hi}}^i]$, a set of
batch sizes $\mathcal{B}=\{16,32,64,128\}$, a round deadline $T$ and a battery
reserve $\mathrm{SoC}_{\min}=0.2$. The deadline is the median, over all
clients, of the nominal training time of the baselines ($E_0=5$, $B_0=128$, no
load). It bounds the stragglers, and thus the idle energy that the other
clients spend waiting for them. The lower bound $E_{\mathrm{lo}}=2$ is the
smallest useful update. The upper bound is
\begin{equation}
  \begin{aligned}
  E_{\mathrm{hi}}^i = \min\biggl(&2E_0,\;
    \left\lfloor \frac{T}{t_i(1,B_{\mathrm{fast}},w_i)} \right\rfloor,\\
  &\left\lfloor \frac{\widehat{\mathrm{SoC}}_i - \mathrm{SoC}_{\min}
    - \delta^{\mathrm{up}}_i}{k_i\,\delta^{\mathrm{ep}}_i} \right\rfloor \biggr),
  \end{aligned}
  \label{eq:ehi}
\end{equation}
clipped to at least $E_{\mathrm{lo}}$. The three terms are a cap at twice the
standard work, the deadline at the fastest batch size, and the energy left
above the reserve. Here $\delta^{\mathrm{ep}}_i$ and $\delta^{\mathrm{up}}_i$
are linear estimates of the SoC spent by one epoch at $\min\mathcal{B}$ and by
the upload, and $k_i$ is a correction factor learned from the consumption the
client reports. The envelope is deliberately loose: the server knows neither
the current load nor the true battery curve, and being precise is the job of
the client.

\textbf{Proposal.}
Before the model, the server sends each of the top-ranked clients a proposal
with its envelope, but no model. The client reads its state and accepts if it
can run $E_{\mathrm{lo}}$ epochs at some $B\in\mathcal{B}$ without crossing its
reserve. It always replies with its SoC, load and charging status, which become
its new report. A refusing client is replaced by the next one in the ranking,
for at most two waves. Its $\eta_i$ is reset, so it does not return to the top
of the ranking at the next round. The model is sent only to the clients that
accepted, so a refusal never wastes a download.

\textbf{Client decision.}
When the model arrives, the client binds its work (late binding). It reads its
SoC, load and charging status, and evaluates the loss $L$ of the global model
on 256 local samples. Then it computes
\begin{equation}
  \begin{gathered}
  u = \min\!\left(1, \frac{L}{\ln C}\right), \qquad
  h = \frac{\mathrm{SoC}-\mathrm{SoC}_{\min}}{1-\mathrm{SoC}_{\min}},\\
  E^\ast = E_{\mathrm{lo}} + \left\lfloor (E_{\mathrm{hi}}-E_{\mathrm{lo}})\,u\,h \right\rceil,
  \end{gathered}
  \label{eq:decision}
\end{equation}
where $C$ is the number of classes and $\lfloor\cdot\rceil$ rounds to the
nearest integer. $\ln C$ is the loss of a random guess, so $u\approx1$ means
that the global model has learned nothing yet on these data, and $h=1$ while
the phone is charging. Much loss and much charge yield many epochs; little of
either yields few. The client then tries the smallest batch size first and
lowers $E$ from $E^\ast$ towards $E_{\mathrm{lo}}$ until two conditions hold
on the predicted trajectory, not on the current state. The SoC at the end of
the round, upload included, must stay above the reserve, and the training time
under the current load must fit the deadline. Only if not even $E_{\mathrm{lo}}$
fits does it move to the next batch size. If the deadline cannot be met at any
batch size, it runs $E_{\mathrm{lo}}$ at the fastest one and declares it. If not
even $E_{\mathrm{lo}}$ fits the reserve, it withdraws from the round (opt-out).
During training, a guard stops the client if it reaches its energy budget, even
in the middle of an epoch.

\textbf{Aggregation and update.}
The client returns the update together with the chosen $(E,B)$, the number of
SGD steps, its initial and final SoC, and its loss. The server discards updates
outside the envelope, charges the declared work to the device, and averages
the updates weighted by $n_i$, as in FedAvg~\cite{mcmahan2017fedavg}. The final
SoC becomes the new report, and $k_i$ is updated as a moving average of the
ratio between the observed and the linearly predicted drain.

\begin{figure}[t]
\small
\begin{algorithmic}[1]
\REQUIRE last reports, static profiles, round $r$, target size $k_r$
\STATE \textbf{Server:} estimate $\widehat{\mathrm{SoC}}_i$ (Eq.~\ref{eq:estimate}) and rank by $s_i$ (Eq.~\ref{eq:score})
\STATE $S \gets \emptyset$, $W \gets 0$
\WHILE{$|S| < k_r$ \AND $W < W_{\max}$}
  \STATE send proposal $\mathcal{E}_i$ (no model) to the next $k_r-|S|$ ranked clients
  \FOR{each reply of client $i$}
    \STATE update the report of $i$
    \IF{$i$ accepts} \STATE $S \gets S\cup\{i\}$ \ELSE \STATE $\eta_i \gets 0$ \ENDIF
  \ENDFOR
  \STATE $W \gets W+1$
\ENDWHILE
\STATE send the model and $\mathcal{E}_i$ to every $i\in S$
\STATE \textbf{Client $i\in S$:} read SoC, $w$, charging; compute $L$ on 256 samples
\STATE compute $u$, $h$, $E^\ast$ (Eq.~\ref{eq:decision})
\FOR{$B\in\mathcal{B}$ in increasing order}
  \FOR{$E = E^\ast$ \textbf{down to} $E_{\mathrm{lo}}$}
    \IF{predicted SoC $\geq \mathrm{SoC}_{\min}$ \AND predicted time $\leq T$}
      \STATE train $E$ epochs at batch $B$ (with the energy guard) and reply
    \ENDIF
  \ENDFOR
\ENDFOR
\STATE if no $(E,B)$ fits: $E_{\mathrm{lo}}$ at the fastest $B$ if only the deadline fails, otherwise opt out
\STATE \textbf{Server:} average the replies whose $(E,B)\in\mathcal{E}_i$, weighted by $n_i$
\end{algorithmic}
\caption{One round of SAGE-smart. The client stops at the first feasible
$(E,B)$.}
\label{alg:smart}
\end{figure}

% =====================================================================
\subsection{Protocols and Cost of Information}\label{sec:protocols}

The client state can enter the decision in three ways. They differ in what
travels on the network, in who pays for it, and in how fresh the state is when
the work is decided. Table~\ref{tab:protocols} counts them per round for $N_a$
alive clients, $k$ selected, $r$ refusals and a fraction $\rho$ of clients that
do not answer a request (out of coverage, switched off); we use $N_a=100$,
$k=20$ and $\rho=0.1$.

\textbf{P1, model-first negotiation.} The server sends the model and an
envelope to the selected clients, which accept, reduce the work, or refuse, and
report their choice with the update. A refusal wastes a download,
$s\,e_{\mathrm{dl}}+E_{\mathrm{s}}\approx46\,\mathrm{J}$. Replacing the client
requires another download and delays the round by its transfer time.

\textbf{P2, lightweight negotiation (SAGE-smart).} The proposal travels before
the model, so a refusal costs an isolated control session, about 2.1~J.
For a client that accepts, the proposal shares the session of the download
that follows. Each wave delays the round by $t_{\mathrm{c}}\approx0.33\,\mathrm{s}$.
The client still binds its work at training time, as in P1.

\textbf{P3, state broadcast.} At the start of the round, the server asks every
alive client for its state and decides $(E,B)$ itself; the client executes. The
cost is fixed. Every responder that is not selected pays a session,
$((1-\rho)N_a-k)\,E_{\mathrm{s}}\approx150\,\mathrm{J}$ per round (about
10~Wh over 250 rounds), and a timeout (1~s) for the clients that do
not answer extends the round. There are no refusals, but the work is decided
on the state at the broadcast. A change of load or a charger plugged in before
training is not seen.

\begin{table}[t]
\centering
\caption{Information cost of the three protocols per round ($N_a=100$, $k=20$,
$\rho=0.1$, $r$ refusals). Energy of the control messages only; the model
exchange (Eq.~\ref{eq:comm}) is the same for all.}
\label{tab:protocols}
\footnotesize
\setlength{\tabcolsep}{3pt}
\begin{tabular}{lccc}
\hline
& P1 model-first & P2 lightweight & P3 broadcast \\
\hline
Who decides $(E,B)$     & client          & client          & server \\
State used              & at training     & at training     & at broadcast \\
Messages before model   & 0               & $k+r$           & $N_a$ \\
Energy per refusal      & 46~J      & 2.1~J     & -- \\
Fixed energy per round  & 0               & $\approx 0$     & 150~J \\
Added latency           & 1 download/refusal & 0.33~s/wave & 1.33~s \\
\hline
\end{tabular}
\end{table}

P2 costs less energy than P1 whenever a refusal occurs, by a factor of about
20. P1 costs less than P3 as long as fewer than
$150/46\approx3$ clients per round refuse. P2 costs less than P3 unless about
70 clients refuse in the same round. The broadcast pays a fixed price for a
state that is fresh at the start of the round. The two negotiations pay only
when a client declines, and they use the state at the moment of training. How
often clients decline depends on the load (Sec.~\ref{sec:results}).

In the experiments of Sec.~\ref{sec:results}, the broadcast is enabled for all
the strategies that select on the device state (SAGE, ESCS and SAGE-smart), so
that all of them pay the same cost for their view of the clients; SAGE-smart
runs P2 on top of it. In simulation, the clients have no life of their own
outside the simulator, so the broadcast cannot return anything the simulator
does not already know. We use it to charge its cost, and SAGE-smart ranks with
the estimate of Eq.~\eqref{eq:estimate}.
\rev{[Update if the protocols become a parameter of the experiments.]}
% NOTE: the estimate uses the same battery model as the simulator, so it is
% wrong only because of recharges and load changes; on real phones it would
% also be wrong because of the model. State it as a limitation.
